\origpage{227}{GRAPH MINORS IV}
{\centering
{\Large Graph Minors. IV. Tree-Width and Well-Quasi-Ordering\par}\medskip

Neil Robertson\par
\textit{The Ohio State University, Department of Mathematics, 231 West 18th Ave., Columbus, Ohio 43210}\par
and\par
P. D. Seymour$^*$\par
\textit{Bellcore, 445 South St., Morristown, New Jersey 07960}\par
Received July 27, 1983; revised September 6, 1989\par
}
\begin{quote}
K. Wagner conjectured that if $G_1, G_2, \dots$ is any countable sequence of finite graphs, then there exist $i, j$ with $j > i \ge 1$ such that $G_i$ is isomorphic to a minor of $G_j$. Kruskal proved this when $G_1, G_2, \dots$ are all trees. We prove a strengthening of Kruskal's result---Wagner's conjecture is true for all sequences in which $G_1$ is planar. We hope to show in a future paper that Wagner's conjecture is true in general, and the results of this paper will be needed for that proof.

\end{quote}

\sectionhead{1. \textsc{Introduction}}

By a ``graph'' we shall mean (except when we say otherwise) in this paper a finite, undirected graph which may have loops and multiple edges. A graph is a \textit{minor} of another if the first can be obtained from a subgraph of the second by edge-contraction.

There is an ever-growing collection of excluded minor theorems in graph theory. (By an ``excluded minor theorem,'' we mean a result asserting that a graph has a specified property if and only if it has no minor isomorphic to a member of a constructively characterized set of graphs.) Perhaps the most imposing of these to date is the result of Archdeacon [1] and Glover, Huneke, and Wang [2], that a graph may be embedded in the projective plane if and only if it has no minor isomorphic to one of 35 specified graphs. An interesting aspect of all excluded minor theorems known is that not only is the set of excluded minors constructively characterized, but also%
\origfootnote{Research partially supported by NSF grant number MCS 8103440}
